\documentclass[12pt]{article}

\usepackage[margin=1in]{geometry}
\usepackage[T1]{fontenc}
\usepackage[utf8]{inputenc}
\usepackage{setspace}
\usepackage{parskip}
\usepackage{titlesec}
\usepackage{lmodern}

\setstretch{1.15}

\titleformat{\section}{\large\bfseries}{\thesection.}{0.5em}{}
\titleformat{\subsection}{\normalsize\bfseries}{\thesubsection.}{0.5em}{}

\begin{document}

\begin{center}
    {\LARGE \textbf{20-Minute Talk Script}}\\[0.5em]
    {\large \textbf{Title: Economics Is Everywhere}}
\end{center}

\vspace{1em}

Good evening everyone,

Thank you very much for being here. It is a pleasure to speak with students and parents today.

When people hear the word \textbf{economics}, they often think of money, inflation, business, or maybe even boring graphs from school textbooks. Some people think economics is only for bankers, stock market experts, or government officials. Others think it is all about numbers and has nothing to do with ordinary life.

But the truth is, economics is much broader, much more practical, and much more interesting than that.

Economics is really about \textbf{choices}.

More specifically, economics is the study of how individuals, families, businesses, and governments make choices when they cannot have everything they want. In economics, we call this \textbf{scarcity}. Scarcity simply means that resources are limited. We all have limited time, limited money, limited energy, and limited opportunities. Because of that, we constantly make decisions.

A student may want more sleep, better grades, a social life, a part-time job, and time for hobbies, but there are only 24 hours in a day. A parent may want to save more, spend more on family needs, travel, invest in education, and pay for rising living costs, but income is limited. Governments want better schools, hospitals, roads, and public services, but budgets are limited too.

So economics begins with a very simple but powerful idea: \textbf{we have to make choices because we live in a world of limits}.

That is why economics matters. It is not just about money. It is about how life works.

Today, I want to do four things. First, I want to explain what economics really is in plain language. Second, I want to show how economics affects everyday life. Third, I want to introduce some exciting fields in economics, including behavioural economics and neuroeconomics. And finally, I want to explain why studying economics is valuable for students and why it opens many doors in the job market.

Along the way, I will use real-life examples, talk about why economics is relevant, and hopefully make it enjoyable too.

\section*{What is economics?}

Let us start with the simplest definition.

Economics is the study of how people make choices under scarcity.

That may sound formal, but in everyday language, economics is about answering questions like these:

Why do prices rise?\\
Why is housing so expensive?\\
Why do some jobs pay more than others?\\
Why do people save or spend money?\\
Why do businesses advertise so much?\\
Why do some countries grow quickly while others struggle?\\
Why do people sometimes make bad financial decisions even when they know better?

Economics helps us think clearly about these questions.

One of the most important ideas in economics is \textbf{trade-offs}. Every time we choose one thing, we give up something else. If a student spends three hours scrolling through social media instead of studying, the cost is not just the time spent. The cost is also the better grade or better understanding they might have gained.

This is called \textbf{opportunity cost}, one of the central ideas in economics.

Opportunity cost is the value of the next best alternative you give up when making a choice.

That idea applies everywhere. If a family spends money on one thing, they may have less for something else. If a government spends more in one area, it may have to spend less in another. If a student chooses one university program, they are also choosing not to pursue other options.

Economics gives us a framework to think more carefully about these choices.

\section*{Economics in everyday life}

One reason economics is so interesting is that it is everywhere.

Take something simple like buying coffee. Imagine a coffee shop on campus in the morning. There is a long line. Why? Because many students want coffee at the same time, but the coffee shop can only serve so many people at once. That is scarcity. That is supply and demand. That is economics.

Or think about concert tickets. A popular artist announces a concert, and tickets sell out in minutes. Why? Because demand is very high, but the number of seats is limited. Then people resell tickets at much higher prices. Suddenly, we are talking about markets, scarcity, willingness to pay, and pricing.

Even grocery shopping is economics. Families compare prices, look for discounts, change what they buy when prices rise, and decide what is necessary and what can wait. Inflation becomes very real when you go to the grocery store and notice that the same basket of food now costs much more than before.

Sometimes I think inflation is when you go shopping for three items and somehow leave the store feeling like you just financed a small car.

But that is exactly why economics matters. It helps explain the forces behind what people experience every day.

Housing is another example. Many families and students ask: why is housing so expensive? Economics tells us to look at supply, demand, construction costs, interest rates, population growth, zoning rules, and local market conditions. Economics does not always give one easy answer, but it helps us identify the right questions.

Economics also shows up in personal decisions. Should you save or spend? Should you work more hours or spend more time studying? Should you buy now or wait? Should you rent or try to buy? These are all economic decisions.

So economics is not something far away in government offices or financial markets. It is part of daily life.

\section*{Big issues economics helps us understand}

Economics is also important because it helps us understand major issues shaping society.

One of those is \textbf{inflation}. Inflation means a general rise in prices over time. When inflation is high, money buys less than before. Families feel this when grocery bills rise, rent increases, and transportation becomes more expensive. Students feel it too, especially when food, books, and housing costs increase.

Another issue is \textbf{unemployment}. Economics studies why people lose jobs, why some industries grow while others shrink, and what policies can support employment. Employment matters not only for income but also for stability, confidence, and well-being.

Then there is \textbf{inequality}. Why do some people earn far more than others? What roles do education, skills, experience, discrimination, and technology play? Economics studies all of these factors.

Economics also helps us understand \textbf{public policy}. When governments change taxes, raise interest rates, expand spending, or introduce new rules, economics helps us analyze who benefits, who loses, and what the unintended consequences might be.

That is one of the most useful lessons in economics: do not just look at the immediate effect of a decision. Also ask what happens next.

For example, if a government wants to make housing more affordable, it cannot just look at one policy in isolation. It also has to think about how builders, renters, buyers, and investors will respond. Economics teaches us to think beyond the first reaction and consider the bigger system.

\section*{Economics is about people}

Sometimes people think economics is cold because it uses numbers, graphs, and statistics. But economics is really about people.

It is about families trying to manage a budget.\\
It is about workers looking for stable jobs.\\
It is about students planning a future.\\
It is about communities trying to create opportunity.\\
It is about governments trying to balance different priorities.

Good economics is not just about equations. It is about understanding human choices and improving lives.

That is why economics connects with so many other fields, including psychology, politics, health, environment, and even neuroscience.

\section*{Interesting fields in economics}

Now let me talk about some of the exciting branches of economics, because economics is not just one narrow subject.

\subsection*{Behavioural economics}

This is one of the most popular areas today.

Traditional economics often assumes that people make rational decisions. Behavioural economics says: not always.

It studies how people actually behave, including the mistakes, emotions, habits, and biases that affect our choices.

For example, why do people procrastinate? Why do we spend money on things we do not really need? Why do we choose short-term pleasure over long-term benefit? Why do people keep paying for subscriptions they forgot they had?

Behavioural economics looks at self-control, habits, overconfidence, fear of losses, peer effects, and many other aspects of real human behaviour.

To put it simply, behavioural economics begins with the idea that human beings are not robots. We are sometimes thoughtful, sometimes emotional, and sometimes irrational.

Or as I like to say, behavioural economics is the field that politely explains why people say, ``I will start saving money on Monday,'' and then order takeout on Tuesday.

This field is very useful in designing better public policies, retirement plans, health programs, and consumer protections.

\subsection*{Neuroeconomics}

This is where economics meets neuroscience and psychology.

Neuroeconomics studies what happens in the brain when people make decisions. It asks questions like: how do we respond to risk? Why do rewards feel good? Why do emotions affect our choices? Why do some people make impulsive decisions?

This is a fascinating field because it goes deeper into how human decision-making actually works.

No, it does not mean economists can read minds. Economists still cannot even predict what a group wants for dinner when everyone says, ``Anything is fine.''

But neuroeconomics helps us understand the brain processes behind trust, risk-taking, reward, and decision-making.

\subsection*{Labour economics}

This field studies jobs, wages, unemployment, skills, education, and the labour market.

Why do some jobs pay more than others? What skills are rewarded in the labour market? What happens when new technology changes how firms operate? Why do wage gaps exist?

This area is especially important for students because it connects directly to careers, education, and earning opportunities.

\subsection*{Environmental economics}

This field studies climate change, pollution, energy, sustainability, and natural resources.

How should society respond to pollution? Should governments use carbon taxes or regulations? How do we protect the environment while still supporting economic growth?

These are very important questions today, and economics plays a major role in addressing them.

\subsection*{Health economics}

This field looks at healthcare costs, access to healthcare, public health policy, and the use of limited medical resources.

\subsection*{Development economics}

This field studies poverty, growth, education, institutions, and living standards in poorer regions and countries.

\subsection*{Financial economics}

This area focuses on banking, financial markets, investing, risk, and asset prices.

So when students study economics, they are not entering one fixed area. They are entering a broad field full of different directions and possibilities.

\section*{Challenges in economics}

Now, economics is powerful, but it is not perfect.

Why? Because people are complicated, and the world changes.

Humans do not always act rationally. We are influenced by emotions, habits, social pressure, uncertainty, and limited information. That makes prediction difficult.

Also, economies are always changing. Technology changes. Markets change. Policies change. Global events can disrupt everything very quickly.

Economics often deals with trade-offs rather than perfect solutions. A policy may help one group but create difficulties for another. A decision that improves one outcome may worsen another.

That is why economists sometimes disagree. But that does not make economics weak. It shows that economics deals with important real-world problems where the answers are not always simple.

\section*{Why study economics?}

Now let us come to the practical question: why should students consider studying economics?

One reason is that economics helps students become strong thinkers. It teaches them how to analyze problems, interpret evidence, compare costs and benefits, and make structured arguments.

Another reason is that economics is very practical. It helps students understand the world they live in and make better personal, professional, and civic decisions.

Economics is also flexible. It combines theory with real-world application. It uses both numbers and ideas. Students who like social issues can find meaning in economics, and students who like data and analysis can find strong tools in economics too.

Perhaps most importantly, economics opens many career paths.

Students with training in economics can work in government, banking, finance, consulting, policy analysis, business, education, research, economic development, data analysis, and many other fields.

An economics degree can also prepare students for graduate school, law school, business school, and public policy programs.

Even when the job title is not ``economist,'' employers value the skills economics teaches: critical thinking, problem-solving, quantitative reasoning, communication, and evidence-based decision-making.

Those are highly marketable skills.

\section*{A message to students and parents}

Let me speak directly to the students first.

You do not need to have your entire future figured out right now. What matters is curiosity and a willingness to learn how the world works.

If you are interested in questions like why prices rise, why some people earn more than others, why people make the choices they do, or how policies affect society, economics may be a very good fit for you.

Now to the parents.

Economics is one of those fields that offers both practical and intellectual value. It helps students build strong analytical skills while also preparing them for a wide range of careers. It gives them tools that remain useful whether they work in business, government, education, finance, policy, or many other areas.

In a changing world, the ability to think clearly, analyze evidence, and make informed decisions is extremely valuable.

Economics helps build exactly those abilities.

And perhaps, after taking economics, students may finally understand why parents keep saying, ``We have to stay within the budget.''

That may be one of the most successful outcomes of all.

\section*{Conclusion}

Let me close with this.

Economics is not just about money. It is about choices, trade-offs, incentives, behaviour, and how society works.

It helps us understand everyday life, from grocery bills and housing costs to careers and public policy. It helps us analyze major issues like inflation, unemployment, inequality, and sustainability. And it offers exciting fields like behavioural economics, neuroeconomics, labour economics, environmental economics, and many more.

Economics is relevant because life is full of decisions, and economics gives us tools to think more clearly about them.

For students, it opens doors. For parents, it offers reassurance that these skills have real value. And for everyone, it provides a better way to understand the world.

Once you start seeing the world through an economic lens, you realize that economics is everywhere.

Thank you very much.

\end{document}